In terms of future work, there are several avenues that could be explored in relation to this research. As discussed in \autoref{Discussion}, the model we have created is a deliberate simplification, but it serves as a structured foundation from which future investigation can proceed. We believe the most productive paths forward fall into two broad categories: expanding the model to address its identified limitations, or moving beyond the conceptual framework entirely to begin testing on a real system.

\subsection{Expanding the Model}
\label{ExpandingTheModel}

The most straightforward path forward here is to address the limitations identified in \autoref{Discussion} incrementally, building on the existing OSM and AESM rather than starting from scratch. A number of those limitations point to well-defined extensions that could meaningfully deepen the model's usefulness.

\subsubsection{Time Series Modelling}
\label{TimeSeriesModelling}

As noted in \autoref{TimeSeries}, the model is a static snapshot rather than a dynamic process. Introducing discrete time steps would allow the simulation to capture the temporal dynamics of the arms race described in \autoref{ArmsRaceDynamic}, most critically the question of whether escalating resource expenditure converges to a stable equilibrium or continues until one agent exhausts its available capacity. The inode timestamp mechanism already provides a natural basis for this extension: since every file operation updates the access, modification, and change timestamps of the affected inode, a time-series model could be built by taking periodic snapshots of the state map and computing the difference between them. Selecting an appropriate sampling interval would itself be a research decision, as a very high frequency would impose significant I/O overhead while a low frequency would risk missing short-lived intermediate states.

\subsubsection{Intermediate States and Memory}
\label{IntermediateStatesMemory}

The current model collapses the condition of an agent into a binary alive or dead distinction. As discussed in \autoref{StatesDiscussion}, a richer model would distinguish between partial corruption, process degradation, and filesystem lockout, each of which carries different strategic implications for both the affected agent and its rival. Introducing these intermediate states would allow the simulation to capture outcomes that the current model cannot represent, such as a partially disabled agent that retains enough capability to obstruct its rival but is no longer able to mount an active offence.

Related to this is the absence of any representation of RAM in the current model, which is discussed in \autoref{TemporaryMemory}. An agent that stores working data in memory rather than on disc would be largely invisible to filesystem-level monitoring, since nothing is written to any path that tools like \texttt{ls} or \texttt{find} would surface. Extending the AESM to include a memory layer would require treating the \texttt{/proc} pseudo-filesystem as both a surveillance channel and a data store, since it exposes live process memory maps, open file descriptors, and environment variables without any corresponding disc writes.

\subsubsection{Privilege Level and Threat Vector Analysis}
\label{PrivilegeThreatVector}

As discussed in \autoref{PrivilegeLevel}, the privilege context of the agents is currently unspecified. We should treat the unprivileged and root cases as distinct scenarios in any future iteration of the model, since the set of viable strategies differs substantially between them. An unprivileged agent cannot terminate processes owned by another user, cannot write to system directories, and cannot modify permissions it does not own, which eliminates a significant portion of the offensive strategies identified in the analysis. Modelling each privilege context explicitly would also give us a natural basis for the formal threat vector analysis described in \autoref{ThreatVectorAnalysis}: once we have fixed the agents' capabilities and specified the privilege domain, it becomes feasible to create a structured mapping of each AESM state category against the available attack methods.

\subsubsection{Agent Asymmetry, Multiple Agents, and Starting Scenarios}
\label{AsymmetryMultipleAgentsScenarios}

There are also a few extensions worth considering that concern the initial conditions of the simulation. First, as discussed in \autoref{AgentSymmetry}, the symmetric two-agent assumption is a useful simplification for the initial model but does not reflect likely real-world deployment scenarios. Introducing configurable asymmetry in resources, capabilities, or knowledge would substantially change the dominant strategies on both sides and would allow us to test how robust our conclusions are across a range of agent configurations, rather than treating the balanced case as the only one worth considering.

Second, extending to three or more agents, as discussed in \autoref{MultipleAgents}, introduces the additional complexity of temporary alliances and defection incentives. When a third agent is present, purely adversarial two-player reasoning breaks down: it may become rational for two agents to cooperate against a third before eventually turning on each other. This is where the formal tools of game theory, particularly coalition formation and iterated defection, could be a genuinely useful analytical framework for this class of multi-agent interaction.

Third, the three starting scenarios outlined in \autoref{StartingScenarios}, specifically the defender head start, undetected infiltration, and resource imbalance, could each be implemented as a parametric variant of the existing model. This would allow us to quantify how sensitive the outcomes are to initial conditions, which is something the current conceptual model simply cannot address. In particular, the undetected infiltration scenario, in which the attacker has already established a foothold by the time the defender's monitoring begins, is the most operationally realistic of the three and is probably the most interesting one to model first.

\subsection{Real System Testing}
\label{RealSystemTesting}

The more ambitious path forward would be to actually move beyond the conceptual model and begin testing the derived strategies in a live environment: specifically, a Linux-based virtual machine, which gives us a tangible and controllable environment without the risks of experimenting on production hardware.

The primary advantage of this approach is realism. The conceptual model necessarily makes simplifying assumptions about command execution, resource consumption, and scheduling, and some of those assumptions may not hold when the agents are implemented as actual processes subject to real kernel scheduling, memory limits, and I/O contention. A live simulation would surface the discrepancies between theoretical strategy and practical effectiveness, and in doing so would either validate or challenge the conclusions we drew from the analysis. In particular, it would give us the first empirical test of the resource-efficiency hypothesis we raised in \autoref{ArmsRaceDynamic}: the claim that the most effective strategy is neither the most aggressive nor the most defensive but the most resource-efficient can only be confirmed through repeated trials.

The resource demands of such a simulation are considerable. Training an AI agent through iterative interaction requires running the simulation many times, updating the agent's policy based on each outcome. The arms race dynamic suggests that convergence may require a large number of iterations, particularly if the agents are symmetric and neither has a structural advantage from the outset. Beyond the computational cost of the training runs themselves, the simulation infrastructure would need to be reset between trials to ensure that residual state from a previous run does not contaminate subsequent results.

There is also a legal and ethical dimension to this work that warrants careful consideration before any implementation begins. A simulation in which two agents attempt to terminate each other's processes and corrupt each other's files is, in its essential structure, an automated adversarial attack scenario. Even when contained within a virtual machine, such a simulation produces techniques and potentially trained models that could in principle be repurposed. Any future implementation should involve the relevant institutional ethics review processes and put appropriate controls in place over how the trained models may be accessed or used. The precedent set by the restricted release of Claude Mythos in 2026 \parencite{MythosNature2026} illustrates that this concern is already being taken seriously at the frontier of AI capability research, and it applies with equal force to research simulations that train agents on adversarial behaviour, even at a much smaller scale.

Despite these challenges, the potential value of a real system simulation is significant. It would let us test the strategies identified in \autoref{DefensiveStrategies} and the agent state categories from \autoref{AgentDerivedStates} empirically rather than theoretically, and would give us concrete data about the conditions under which the arms race dynamic stabilises. If we also included the starting scenarios from \autoref{StartingScenarios} alongside the symmetric baseline, it would shed light on just how sensitive the outcomes are to initial conditions, which is something we simply cannot resolve within the conceptual framework alone.